\documentclass{report}

% Packages for images.
\usepackage{graphicx}
\usepackage{float}

% Packages for mathematical formulae/symbols.
\usepackage{amsmath}
\usepackage{amssymb}
\usepackage{amsfonts}

% Packages for code snippets and algorithms.
\usepackage{listings}
\usepackage{xcolor}
\definecolor{codegreen}{RGB}{0,128,0}
\definecolor{codegray}{RGB}{110,110,110}
\definecolor{codepurple}{RGB}{163,21,21}
\definecolor{codeblue}{RGB}{0,0,180}
\definecolor{backwhite}{RGB}{255,255,255}
\definecolor{backcolour}{RGB}{248,248,248}
\lstdefinestyle{mypython}{
    language=Python,
    backgroundcolor=\color{backwhite},
    basicstyle=\ttfamily\small,
    keywordstyle=\color{codeblue}\bfseries,
    commentstyle=\color{codegreen}\itshape,
    stringstyle=\color{codepurple},
    numberstyle=\tiny\color{codegray},
    numbers=left,
    numbersep=8pt,
    frame=single,
    rulecolor=\color{black!20},
    breaklines=true,
    showstringspaces=false,
    tabsize=4,
    keepspaces=true
}
\lstset{style=mypython}
\usepackage[ruled,vlined]{algorithm2e}

% Creating the Example boxes.
\definecolor{exampleback}{RGB}{240, 248, 240}
\definecolor{exampleframe}{RGB}{34, 112, 62}
\usepackage[most]{tcolorbox}
\newtcolorbox{examplebox}[1][]{
    colback=exampleback,
    colframe=exampleframe,
    coltitle=white,
    fonttitle=\bfseries,
    title={Example\ifblank{#1}{}{: #1}},
    arc=2pt,
    boxrule=1pt
}

% Packages for styling and layout.
\usepackage{booktabs}
\usepackage{array}
\usepackage{geometry}
\usepackage{hyperref}

% Packages for plotting graphs.
\usepackage{tikz}
\usepackage{pgfplots}
\pgfplotsset{compat=1.18}

% Packages for compilation.
\usepackage{microtype}

\begin{document}

\begin{titlepage}
    \centering
    \vspace*{2cm}
    
    % Subject Title.
    {\Huge \bfseries Fundamentals of Data Science \par}
    \vspace{1.5cm}
    
    % Master's Degree and Faculty.
    {\Large Master's Degree in Data Science, Sapienza Università di Roma \par}
    {\large Faculty of Information Engineering, Informatics and Statistics \par}
    \vspace{0.5cm}
    
    % Department.
    {\large Department of Computer Science \par}
    \vspace{3cm}
    
    % Author.
    {\Large \textbf{Author:} Gianmaria Romano \par}
    \vspace{3cm}

    % Reference Bibliography
    {\Large Reference Bibliography \par}
    {\large M. P. Deisenroth, A. A. Faisal, C. S. Ong: \href{https://mml-book.github.io/}{\textit{Mathematics for Machine Learning}} \\
    S. Prince: \href{https://udlbook.github.io/udlbook/}{\textit{Understanding Deep Learning}} \\
    J. Watt, R. Borhani, A. K. Katsaggelos: \href{https://github.com/neonwatty/machine-learning-refined/tree/main}{\textit{Machine Learning Refined}} \\}
    \vfill
    
    % Academic Year.
    {\large Academic Year 2026/2027 \par}
    \vspace*{1cm}
\end{titlepage}

\clearpage

%\maketitle
\begin{tcolorbox}[colback=gray!5!white, colframe=gray!75!black, title=Disclaimer]
    This document contains notes for the \textbf{Fundamentals of Data Science} course, given by \textbf{Professors Cinelli and Spinelli} as part of the Master’s Degree in Data Science at \textbf{Sapienza Università di Roma} during the \textbf{first} semester of Academic Year \textbf{2026/2027}.
    
    \smallskip
    \textbf{Note: }These notes may not cover all details of each lecture, so it is highly suggested to consult the Professor's official resources and materials. 
    
    \smallskip
    \textbf{Important: }These notes are free to share and use, but please remember to credit me as the author and do not try to hide/remove this page.
\end{tcolorbox}

\begin{tcolorbox}[colback=gray!5!white, colframe=gray!75!black, title=Contacts]
    For further information or requests, feel free to \href{https://t.me/toritosenso}{contact me privately} or reach out through my repository: \href{https://github.com/GianmariaRomano/Data-Science-Notes}{https://github.com/GianmariaRomano/Data-Science-Notes}.
\end{tcolorbox}

\begin{tcolorbox}[colback=gray!5!white, colframe=gray!75!black, title=License]
    These documents are distributed under the \href{https://www.gnu.org/licenses/gpl-3.0.en.html}{GNU General Public License 3.0}, a form of copyleft intended for use on a manual, textbook or other documents. Material licensed under the current version of the license can be used for any purpose, as long as the use meets certain conditions:
    \begin{itemize}
        \item All previous authors of the work must be attributed.
        \item All changes to the work must be logged.
        \item All derivative works must be licensed under the same license.
        \item The full text of the license, unmodified invariant sections as defined by the author if any, and any other added warranty disclaimers (such as a general disclaimer alerting readers that the document may not be accurate for example) and copyright notices from previous versions must be maintained.
        \item Technical measures such as DRM may not be used to control or obstruct distribution or editing of the document.
    \end{itemize}
\end{tcolorbox}

\tableofcontents

\chapter{Introduction to data science}
\section{The lifecycle of data science}
Data science can be seen as an interdisciplinary field that focuses on extracting actionable insights and knowledge from structured and unstructured data using techniques based on statistical and mathematical methods, as well as computer science algorithms and domain-specific knowledge. \newline
Generally speaking, a data science pipeline starts with understanding the business objectives and translating them into machine learning objectives. \newline
Next, after collecting and cleaning data, the pipeline focuses on selecting and training a model and evaluating the results for deployment, ultimately monitoring model efficacy and restarting the pipeline, either by collecting or data or re-training the model, whenever needed.
\section{The role of machine learning in data science}
Artificial intelligence is a sub-field of computer science that focuses on creating agents that can perform tasks requiring human intelligence and problem solving. \newline
More specifically, Tome Mitchell defines machine learning as ``the study of algorithms that improve their performance $\displaystyle P$ at some task $\displaystyle T$ with experience $\displaystyle E$``, typically focusing on tasks that would be too complex to just hard-code.
\begin{examplebox}[\textbf{Example: }Playing checkers]
    The task $\displaystyle \langle P, T, E \rangle$ to train an agent to play checker can be defined using the following components:
    \begin{itemize}
        \item The agent's performance $\displaystyle P$ is evaluated as the proportion of matches it wins against an arbitrary opponent.
        \item The task $\displaystyle T$ to perform is to train the model to make the correct moves during a match.
        \item The experience $\displaystyle E$ used to train the agent comes from matches the agent plays against itself.
    \end{itemize}
\end{examplebox}
\noindent While machine learning and statistics typically use the same instruments and algorithms to discover patterns within data, the two fields are normally considered to be different areas because statistics follows a more rigorous inference-based approach, while machine learning focuses on making predictions using scalable and autonomous techniques.

\end{document}